\documentclass[tikz, border=2pt]{standalone}
\usepackage{ctex}
\usetikzlibrary{positioning, calc, arrows.meta, shapes.geometric, fit, backgrounds}

\begin{document}
\begin{tikzpicture}[
    font=\sffamily\small,
    component/.style={draw, rounded corners, minimum width=3cm, minimum height=1.5cm, align=center, fill=white, thick},
    hardware/.style={draw, trapezium, trapezium left angle=70, trapezium right angle=110, minimum width=3cm, minimum height=1.2cm, align=center, fill=gray!10, thick},
    arrow/.style={->, >=Latex, thick},
    label/.style={font=\footnotesize}
]

% X Client
\node[component] (client) {X Client\\(应用程序)};

% X Server
\node[component, right=4cm of client] (server) {X Server\\(显示服务器)};

% Hardware
\node[hardware, right=1.5cm of server] (hw) {硬件\\(键盘, 鼠标, 屏幕)};

% Communication between Client and Server
\draw[arrow] ([yshift=0.3cm]client.east) -- node[above, label] {请求 (Requests)} ([yshift=0.3cm]server.west);
\draw[arrow] ([yshift=-0.3cm]server.west) -- node[below, label] {事件 (Events) / 回复 (Replies)} ([yshift=-0.3cm]client.east);

% Label for the connection
\node[label, above=0.8cm] at ($(client)!0.5!(server)$) {X 协议 (X Protocol)};
\node[label, below=0.8cm] at ($(client)!0.5!(server)$) {网络 / 本地套接字};

% Communication between Server and Hardware
\draw[<->, thick] (server) -- (hw);

% Grouping for context (optional, but helps visualize separation)
% \node[draw, dashed, fit=(client), label=above:机器 A] {};
% \node[draw, dashed, fit=(server)(hw), label=above:机器 B] {};

\end{tikzpicture}
\end{document}
